An app or agent that acts for a Soroban smart account without a signature prompt
for each call needs a bound on what it may do: which contracts and functions,
which recipients, how much over time, and until when. Two designs build that
bound on OpenZeppelin's account model~\cite{oz-stellar-accounts}.
oz-policy-builder~\cite{oz-policy-builder} attaches a shared predicate
interpreter to the account's context rules as a policy, beside OpenZeppelin's own
policies, and writes its predicates with an off-chain synthesiser. Perch links
one evaluator into the account and decides every grant there. The Perch paper,
\perchdoc{main}, gives the comparison one section; this companion gives it in
full, with sources, and ends with what Perch should learn.

\paragraph{Sources and runs.} We compare Perch~\cite{perch-src} at \fact{perch-pin-short}, whose
evaluator is rule-core~\cite{rule-core} at \fact{rule-core-pin-short};
oz-policy-builder at \fact{opb-pin-short} (\fact{opb-pin-date}), release
\fact{opb-release} and \fact{opb-commits-after-release} later commits, none under
its \srcpath{contracts/} directory; and OpenZeppelin \code{stellar-contracts}
\fact{oz-v080-tag} (\fact{oz-v080-commit}). Claims about oz-policy-builder are
read from its public documentation and source at that pin, cited by path and
line (unqualified paths are oz-policy-builder's), or come from our runs. We built
it, and its \fact{opb-rust-test-runs} Rust test runs and \fact{opb-ts-tests}
TypeScript tests passed; we ran its record, synthesise, verify and install flow
against an OpenZeppelin \fact{oz-v080-tag} account on
\fact{opb-local-network}, and the calls outside the rule were refused with the
codes its documentation gives; and we ran one task through both designs in
soroban-sim~\cite{soroban-sim} (\S\ref{sec:same-task}).

\paragraph{Disclosure.} The authors designed Perch and work at The Aha Company.
soroban-sim, the host in which both designs ran, is published at
\code{\fact{soroban-sim-repo}} (\fact{soroban-sim-license}).

\paragraph{Shared ground.} Both designs use OpenZeppelin's context-rule model: an
authorization entry names, in its \code{context\_rule\_ids}, one rule for each
context (each call or contract creation the entry authorizes), and that rule
alone decides the context. Perch reproduces OpenZeppelin's authorization types
and its two signer kinds, \emph{external} (a key a verifier contract checks) and
\emph{delegated} (an address whose own authorization the host checks), but
attaches no policies. Both judge each context alone, so a call that authorizes
two contexts needs a rule for each (\srcpath{docs/architecture.md:307-309}). Both
deny by default.

\paragraph{The trade-off.} OpenZeppelin with oz-policy-builder composes a
decision from contracts that can be swapped and audited separately; its authors
call the interpreter ``one audit-once engine'' to which ``every policy is
declarative data'' (\srcpath{README.md:21-24}). Perch
puts one grant document behind one decision point: the account's linked
evaluator checks target, functions, arguments, threshold, expiry and window cap
together, against state held only in the account. Composition buys flexibility
and reuse; one decision point gives one place to check and to prove. Eight words
mean different things in the two systems (Table~\ref{tab:tab-terms}); we use them
in Perch's sense unless a sentence says otherwise.

\perchtable{tab-terms}
